\section{Replication Implementation}

\subsection{Repository structure}
\begin{itemize}
  \item The replication is implemented as a configurable Python project with source code under \artifact{src/}, experiment definitions under \artifact{config/}, and paper-level scripts under \artifact{scripts/}.
  \item The original paper source is preserved under \artifact{docs/papers/arxiv-1612.03770v2/}, which allows the report to compare local implementation choices against the paper text and figures.
  \item The dynamic-capacity literature overview is stored as \artifact{docs/papers/dynamic_capacity_taxonomy_original_figures_cl.html}; it should be treated as a local taxonomy draft, not as a citable scholarly source.
\end{itemize}

\subsection{Implemented experiment families}
\begin{itemize}
  \item \artifact{scripts/run_paper_experiments.py} defines a compact paper suite covering MNIST classical learning, MNIST classical learning with intrinsic replay, MNIST NDL, MNIST NDL with intrinsic replay, and SD-19 NDL variants.
  \item \artifact{scripts/run_paper_config.py} runs YAML-defined paper experiments from \artifact{config/paper/}, supports repeated runs, assigns seeds, logs to MLflow when available, and writes summary artifacts.
  \item The key paper configurations are \artifact{mnist_cl}, \artifact{mnist_cl_ir}, \artifact{mnist_ndl}, \artifact{mnist_ndl_ir}, \artifact{sd19_ndl}, \artifact{sd19_ndl_ir}, and \artifact{figure4}.
\end{itemize}

\subsection{Implementation claims to document}
\begin{itemize}
  \item Classical learning baselines isolate the effect of ordinary parameter updates without architecture growth.
  \item Intrinsic replay baselines isolate the effect of replay without architecture growth.
  \item NDL runs isolate the effect of neurogenesis-style hidden-layer growth.
  \item NDL with replay evaluates whether growth and replay act as complementary mechanisms.
  \item SD-19 variants test whether the implementation generalizes beyond the smaller MNIST setting.
\end{itemize}

\subsection{Implementation evidence to include}
\begin{itemize}
  \item Record the exact command used for each run, including the config file path and seed policy.
  \item Record the final hidden sizes and total parameter counts because architecture growth is the central treatment.
  \item Record update counts for pretraining, neurogenesis, incremental training, and the full run when the summary JSON exposes them.
  \item Record runtime and hardware environment as reproducibility metadata rather than as the main result.
\end{itemize}
